% Proposed consequence of audited EL and FF.4; pending independent review.
\paragraph{An exceptional perfect-$K_0$ presentation of the ordinary
top quotient (proposed comparison).}
Retain all hypotheses of the general endpoint theorem EL: $k$ is
algebraically closed of characteristic $p>0$; $R$ is an excellent
henselian local coefficient-$k$ ring of algebraizable dimension three;
$Z\to\operatorname{Spec}R$ is a regular projective modification,
an isomorphism away from its closed point; its reduced closed fibre
$E$ is a simple normal crossing divisor; and an effective Cartier
divisor $D$ with full support $E$ has ample conormal
$\mathcal O_D(-D)$.  Put $X=Z\setminus E$.

Define
\[
 S(E)=F_\gamma^2K_0(E)_{\mathbf Q}.
\]
Let $J_K(E)\subset S(E)$ be generated by all differences
\begin{equation}
          u^*\beta-v^*\beta,
 \qquad \beta\in F_\gamma^2K_0(E\times P)_{\mathbf Q},
 \quad u,v\in P(k),                              \tag{EP.1}
\end{equation}
where $P/k$ runs over smooth connected projective curves and the
pullbacks are along the actual two product sections.  This definition
uses perfect $K$-theory on the possibly singular projective product.
It does not replace it by coherent $G$-theory or by cycles on the
individual normal-crossing strata.

Restriction and the rational Chern character give an isomorphism
\begin{equation}
 \rho:\mathrm{CH}^2(Z)_{\mathbf Q}
       \xrightarrow{\sim}S(E).
                                                        \tag{EP.2}
\end{equation}
Let
\[
 \theta:\mathrm{CH}_1(E)_{\mathbf Q}\longrightarrow S(E)
\]
be the composite of proper cycle pushforward into $Z$ with $\rho$.
The comparison to be proved is
\begin{align}
 \rho\bigl(\mathrm{CH}^2(Z)_{\mathrm{alg},\mathbf Q}\bigr)
       &=J_K(E),                                      \tag{EP.3}\\
 B^2(X)_{\mathbf Q}
       &\simeq S(E)\big/
          \bigl(J_K(E)+\theta\mathrm{CH}_1(E)_{\mathbf Q}\bigr).
                                                        \tag{EP.4}
\end{align}
Here ordinary algebraic equivalence means actual fixed-product
families.  In particular the subgroup on the right of (EP.3) no
longer involves a lift of a parameter family to the resolution.
The exceptional pushforward $\theta$ still records the given
embedding $E\subset Z$, as it should.

\emph{1. The point restriction and the weight-two identification.}
FF.4 gives
$K_0(Z)_{\mathbf Q}\simeq K_0(E)_{\mathbf Q}$ by actual restriction.
Its proof uses KST and levelwise nilpotent invariance, together with
$K_{-1}(R)=0$ for the henselian local ring, by
\citet[Theorem 3.7 and Remark (ii), p.~8]{frDrinfeld2004Tate}.
Restriction commutes with lambda and Adams operations.

Since $E$ is a projective surface, $F_\gamma^3K_0(E)=0$.
The Adams action on $F_\gamma^i/F_\gamma^{i+1}$ is $2^i$ for
$\psi^2$\citep[Chapter II, Proposition 4.9 and Remark 8.8.5]
{D_Weibel2013KbookII}.  Over $\mathbf Q$ the three distinct roots
$1,2,4$ give the eigenspace decomposition.  Explicitly
$(\psi^2-1)(\psi^2-2)/6$ is the projector onto $F_\gamma^2$.
Thus $S(E)$ is exactly the Adams-weight-two subgroup.

For $Z$, the rational Chern character identifies the weight-two
subgroup with $\mathrm{CH}^2(Z)_{\mathbf Q}$.  The strict-local
use is justified as in TC: the projective resolution and every
finite collection of bundles, cycles, relations and Chern
operations descend to one pointed etale finite-type model.
After shrinking around its proper exceptional fibre the model
is regular, hence smooth over the perfect field $k$.
The smooth rational Chern-character isomorphism
\citep[Chapter II, Corollary 8.9.1, p.~84]{D_Weibel2013KbookII}
passes to this filtered limit.  Restriction therefore gives (EP.2).
The same reasoning applies to $Z\times P$ and its finite cycle
families; it does not require that the singular base be smooth.

\emph{2. Every exceptional perfect-family relation lifts at its endpoints.}
The map
\[
       K_0(D\times P)_{\mathbf Q}
             \longrightarrow K_0(E\times P)_{\mathbf Q}
\]
is an isomorphism by nilpotent invariance in characteristic $p$
\citep[Theorem A]{rootGeisserHesselholt2011}.  It is a lambda-ring
map by perfect pullback.  Its inverse preserves rank and the
rational gamma filtration: one may either use the inverse
lambda-ring map, or lift each rank-zero generator in an expression
for $F_\gamma^2$ and apply its gamma operations.  Rational gamma
operations are determined by the universal special lambda-ring
identities after clearing denominators.  This verifies the
filtration assertion without identifying ordinary cycles on the
nonreduced scheme.

Given $\beta$ in (EP.1), let $\widetilde\beta$ be its preimage in
$F_\gamma^2K_0(D\times P)_{\mathbf Q}$.  EL constructs an actual
$b\in K_0(Z\times P)_{\mathbf Q}$ whose two endpoint restrictions
to $E$ are precisely those of $\beta$.  Since these endpoints lie
in $S(E)$, the point isomorphism in Step~1 forces the two endpoint
classes of $b$ on $Z$ to be of weight two.  The actual cycle
$\operatorname{ch}_2(b)\in\mathrm{CH}^2(Z\times P)_{\mathbf Q}$
therefore has the prescribed endpoint difference under $\rho$.
It is an ordinary algebraic-equivalence relation.  This proves
$J_K(E)\subset\rho(\mathrm{CH}^2(Z)_{\mathrm{alg},\mathbf Q})$.
The full exceptional class $\beta$ need not lift on the whole
product: only its endpoints were used.  No assertion about
$K_{-1}(\operatorname{Spec}R\times P)$ is needed.

\emph{3. Every ordinary cycle relation restricts to such a perfect family.}
Conversely consider a relation given by a cycle family
$z\in\mathrm{CH}^2(Z\times P)_{\mathbf Q}$ and two fibres.
On a finite-type smooth model, take the inverse rational Chern
character of the element having codimension-two component $z$
and all other components zero.  Pull it to $Z\times P$.
This is an actual class
$b\in F_\gamma^2K_0(Z\times P)_{\mathbf Q}$ with
$\operatorname{ch}_2(b)=z$; it has pure weight two.
Restriction to $E\times P$ preserves $F_\gamma^2$, so its
endpoint difference belongs to $J_K(E)$.  Chern character and
pullback commute, which proves the reverse inclusion in (EP.3).

It suffices here to use projective parameter curves.  Any ordinary
relation originally given over a smooth nonproper curve can be
compactified: extend that curve smoothly to a projective curve,
and close the finitely many horizontal integral cycle components
in the product.  They remain flat over the regular parameter
curve near the two chosen points, and their refined fibres there
are unchanged.  This closure operation is made on a single finite
type model if necessary and then pulled to the strict-local
scheme.  No henselized parameter product is substituted.

\emph{4. Localization preserves the exact ordinary relation subgroup.}
Cycle localization gives
\[
 \mathrm{CH}_1(E)_{\mathbf Q}\longrightarrow
 \mathrm{CH}^2(Z)_{\mathbf Q}\longrightarrow
 \mathrm{CH}^2(X)_{\mathbf Q}\longrightarrow0.
\]
Moreover the image of
$\mathrm{CH}^2(Z)_{\mathrm{alg},\mathbf Q}$ is exactly
$\mathrm{CH}^2(X)_{\mathrm{alg},\mathbf Q}$.
One inclusion is restriction of actual families.  For the other,
close an actual family on $X\times P$ in $Z\times P$; refined
pullback along either parameter point then restricts to the
original fibre on $X$.  This can again be performed on one smooth
finite-type model with finitely many cycle components.  Thus each
ordinary relation on $X$ is the restriction of an ordinary relation
on $Z$, up to the allowed exceptional cycles.  Quotienting the
localization sequence by these subgroups, and using (EP.2)--(EP.3),
gives (EP.4).

\emph{Scope of the reduction.}
The nontrivial input is EL: (EP.3) proves that every actual proper
perfect-family relation on the exceptional surface is realized
by an ordinary cycle relation on the regular resolution, even
when the original family itself has a nonzero negative-$K$
boundary.  This removes the ambient endpoint-lifting obstruction
under the displayed three-dimensional resolution hypotheses.
It does not compute $J_K(E)$ for an arbitrary normal-crossing
surface.  In particular families assembled separately on double
curves can have nonzero restrictions at triple points and need
not glue to a class on $E\times P$; those compatibility conditions
are still part of (EP.1).  The comparison is therefore stronger
than retaining the original liftability condition, but does not
by itself prove universal vanishing or finiteness of the right
side of (EP.4).
